\section{Lattice normalization and a validated plaquette example}\label{sec:applications}

Let $T^a=\sigma^a/2$ and write a plaquette holonomy as
\[
 U_P=\exp(iX_P^aT^a),\qquad
 \frac12\operatorname{Tr}U_P=\cos\theta_P,\qquad
 \theta_P=\frac{|X_P|}{2}\quad (|X_P|\le2\pi).
\]
The one-plaquette heat-kernel action is $s_t(\theta_P)$.  In spin notation
$n=2j$, so the character weight in \eqref{eq:character} is
$e^{-2t\,j(j+1)}$.  A convention that writes the weight as
$e^{-\tau C_2(j)}$ therefore has $\tau=2t$.

For a smooth continuum field,
\[
 X_P^a=a^2g_0F_{\mu\nu}^a(x)+O(a^3),
\]
in the convention where $g_0$ appears in the plaquette exponential.  Hence
\begin{equation}\label{eq:weak-field-lattice}
 s_t(U_P)=\frac{a^4g_0^2}{8t}F_{\mu\nu}^aF_{\mu\nu}^a
+O\!\left(\frac{a^5}{t}\right)+O(a^4g_0^2|F|^2),
\end{equation}
up to the field-independent term $-\log C_t$.  The final matching of $t$ to
the conventional bare coupling depends on whether $g_0$ is included in
$F_{\mu\nu}$ and on the normalization of the continuum action; equation
\eqref{eq:weak-field-lattice} makes that conversion explicit rather than
silently identifying $t$ with a standard $\beta$.

\begin{corollary}[Local nonlinear budget]\label{cor:lattice}
Fix $M\ge3$ and an identity chart satisfying
Theorem~\ref{thm:action}.  Then
\begin{equation}\label{eq:scaledbudget}
 \left|\partial_\theta^m\left(ts_t-\frac{\theta^2}{2}\right)\right|
 \le t(G_m+B_m),\qquad 1\le m\le M.
\end{equation}
For a one-parameter link variation with class angle $\theta_P(x)$,
\begin{align}
 \frac{\dd^2}{\dd x^2}s_t(\theta_P(x))
 &=s_t''(\theta_P)(\theta_P')^2+s_t'(\theta_P)\theta_P'',\label{eq:link-chain}\\
 \left|\frac{\dd^2}{\dd x^2}
 \left(s_t-\frac{\theta_P^2}{2t}\right)\right|
 &\le R_2(\theta_P')^2+R_2|\theta_P|\,|\theta_P''|,\label{eq:link-chain-bound}
\end{align}
as long as the path remains in the chart.  Radial plaquette variations have
$\theta_P'=1$ and $\theta_P''=0$.
\end{corollary}

For the rationally validated instance $t_0=1/5$ and $\rho=1$, the record gives
$R_2=G_2+B_2<0.455$.  Therefore every $0<t\le1/5$ satisfies
\begin{equation}\label{eq:numeric-convexity}
 \lambda_{\rm rad}(s_t)>4.545,\qquad
 \lambda_{\rm tan}(s_t)>2.75
 \quad\text{on }0\le\theta\le1.
\end{equation}
At $t=1/5$, the scaled-action derivative corrections obey
\[
 tR_2<0.091,\qquad
 t(G_3+B_3)<0.058,\qquad
 t(G_4+B_4)<0.101.
\]
Thus a commuting radial plaquette path
$U_P(x)=\exp(2i(\theta+x)T^3)$ has a validated second derivative at least
$4.545$, while its third and fourth scaled derivatives are bounded by the
last two displayed constants.  This is a concrete finite-plaquette
coercivity and nonlinearity statement, not a statement about a physical mass.

For a full lattice, composition with link coordinates requires bounds on the
first and higher derivatives of the noncommutative plaquette product, and
summation over plaquettes requires a specified operator norm.  Gauge
directions also prevent one-plaquette convexity from becoming a global
strict-convexity statement on link space without gauge reduction.

Likewise, Theorem~\ref{thm:tail} concerns a random group element with density
$p_t$ relative to normalized Haar measure.  In a four-dimensional lattice
measure, plaquettes share links and obey nonabelian constraints.  The theorem
does not determine the interacting plaquette marginal or provide the
large-field estimate needed for a renormalization-group argument.
